\chapter{Compressed packet simulation}
\label{ch:compressed-packets}
\index{compressed representation}

\section{Beyond packet dimension}
The small-packet theorem stores one amplitude per packet-basis state. This is not the only classical representation. A state in an exponentially large packet may lie on a polynomially parameterized manifold. Stabilizer states, Gaussian states, matrix-product states of bounded bond dimension, and low-rank operator states are standard examples of compact internal encodings.

Let $\mathcal R_n$ be a class of classical descriptions, and let
\[
\operatorname{decode}:\mathcal R_n\to P_n\Hh_n
\]
associate a quantum state with each description.

\begin{theorem}[Compressed-packet simulation]
\label{thm:compressed-packet}
A polynomial-time uniform quantum circuit family lies in $P$ if there is a polynomial-size classical representation satisfying:

\begin{enumerate}
\item the initial state has an efficiently constructible description $r_0\in\mathcal R_n$;
\item every gate induces a polynomial-time update $r\mapsto T_j(r)$ with
\[
\operatorname{decode}(T_j(r))
=
G_j\operatorname{decode}(r);
\]
\item all intermediate description sizes remain polynomial;
\item the acceptance probability is computable from the final description in polynomial time.
\end{enumerate}
\end{theorem}

\begin{proof}
Iterate the polynomial-time classical updates on the polynomial-size descriptions and evaluate the final measurement. The algorithm never reconstructs the full state vector.
\end{proof}

\section{Two-axis complexity}
The records suggest the pair
\[
(D_{\rm packet},\chi_{\rm internal}),
\]
where $D_{\rm packet}$ measures active packet size or count and $\chi_{\rm internal}$ measures the representation cost inside the packet. Efficient simulation may follow from either coordinate:

\begin{itemize}
\item fixed weight has polynomial $D_{\rm packet}$ and can use explicit amplitudes;
\item stabilizer states may have exponential support but polynomial $\chi_{\rm internal}$;
\item Gaussian states may occupy large number sectors but have polynomial covariance descriptions;
\item half-filled generic states may have both coordinates exponential.
\end{itemize}

\begin{figure}[H]
\centering
\includegraphics[width=.72\textwidth]{two_axis_complexity.pdf}
\caption{Packet dimension and internal representation cost are independent classical resources.}
\end{figure}

\section{Hybrid transfer graphs}
Suppose packets are nodes of a sparse graph and each node carries an internal representation. A gate may update a node internally or transfer a compact component to neighboring nodes. The simulation cost depends on
\[
\sum_{\alpha\in A_t}\operatorname{size}(r_\alpha)
+
\#\{\text{nonzero transfers at time }t\}.
\]
This model combines symmetry selection rules with compact state classes.

\section{Closure is the decisive condition}
A compact representation may be destroyed by one gate. Clifford tableaux fail under generic non-Clifford rotations. Gaussian covariance matrices fail under generic non-Gaussian operations. Low-bond tensor networks may acquire exponential bond dimension. The theorem therefore requires closure or a proved polynomial growth bound, not merely a compact initial state.

\begin{principle}[Representation growth audit]
The correct complexity variable is the maximum representation size reached along the computation, not the size of the initial or final state alone.
\end{principle}
